\myChapter{Revue de la litterature sur la generation d'applications web pilotees par des modeles BPMN}{Revue de la litterature}
\label{chap:revue-litterature-prooweb}
\myMiniToc{mysection}{Sommaire}

\mySection{Introduction}{}

La transformation de modèles de processus métier en applications web constitue un défi scientifique et technique qui mobilise plusieurs champs disciplinaires : la modélisation par BPMN, les systèmes de gestion de processus, l'ingénierie dirigée par les modèles et les plateformes low-code. Si chacun de ces domaines a fait l'objet de contributions significatives, les travaux qui les articulent de manière cohérente en vue de produire des applications web complètes, maintenables et dont le code source reste maîtrisable, demeurent rares et fragmentés.

Ce chapitre propose une revue structurée de la littérature visant à situer ProOWeb dans ce paysage scientifique. La démarche, conduite sur un corpus de travaux publiés entre 2015 et 2026, s'inspire des principes de rigueur et de transparence recommandés par PRISMA 2020, adaptés au contexte de l'informatique et du génie logiciel. Elle est guidée par une question centrale : quelles approches permettent de passer de modèles BPMN vers des artefacts logiciels web, et quelles lacunes justifient la conception d'un framework de génération tel que ProOWeb ?

L'analyse est conduite selon quatre axes. Le premier porte sur BPMN en tant que langage de modélisation et support à l'automatisation des systèmes orientés processus. Le deuxième examine les BPMS classiques, leurs capacités de déploiement applicatif et leurs limites structurelles face aux exigences d'une application web complète. Le troisième traite de l'ingénierie dirigée par les modèles et des mécanismes de génération de code, en interrogeant leur aptitude à produire des artefacts cohérents, évolutifs et intégrables. Le quatrième confronte les approches model-driven aux plateformes low-code et no-code, dont l'essor récent soulève des questions de gouvernance, de maintenabilité et de dépendance technologique.

\mySection{BPMN dans la conception et l'automatisation des systèmes orientes processus}{}

BPMN occupe une position centrale dans les systèmes d'information orientes processus. Le BPM, tel qu'il est présente par van der Aalst \cite{aalst2013}, vise à concevoir, exécuter, surveiller et améliorer les processus organisationnels. Dans cette perspective, le modèle de processus n'est pas seulement un document graphique. Il sert de support à l'analyse, à la communication et parfois à l'automatisation.

Dans les travaux de Drave et al. \cite{drave2022}, le processus est traite comme un artefact d’ingénierie logicielle. Les auteurs montrent qu'un système d'information oriente processus doit combiner plusieurs dimensions : les données, les activités humaines, les activités automatiques et les interactions avec les utilisateurs. Cette vision est importante pour ProOWeb, car elle rappelle que BPMN ne peut pas être considéré isolement. Le modèle de processus doit être relie à des structures de données, des écrans, des services et des règles applicatives.

Les travaux de transformation BPMN vers IFML ou vers des interfaces graphiques confirment ce rôle de BPMN comme point d'entrée. Sajji et al. proposent une démarche de transformation de BPMN vers IFML dans une logique MDA \cite{sajji2022}, puis prolongent cette perspective vers la génération d'interfaces utilisateur \cite{sajji2023}. Laaz et al. \cite{laaz2020} montrent également que BPMN peut être combine à des ontologies de domaine au niveau CIM afin d'enrichir la transformation vers des modèles d'interaction. Dans un autre contexte, Fedeli et al. \cite{fedeli2024} montrent que BPMN peut aussi être enrichi pour prendre en compte des processus métier connectes à des environnements techniques spécifiques.

Cependant, BPMN reste principalement un langage de contrôle du flux. Il indique qu'une activité existe, qu'une décision doit être prise ou qu'un chemin est possible, mais il ne décrit pas complètement les champs d'un formulaire, les droits d’accès, la structure des données, les validations ou les composants logiciels à produire. Pour ProOWeb, la conclusion est claire : BPMN doit être conserve comme point d’entrée métier, mais il doit être complété par des spécifications supplémentaires.

\mySection{BPMS classiques et applications web métier}{}

Les BPMS classiques apportent une réponse concrète au besoin d’exécuter des processus. Ils gèrent les instances, les taches, les transitions, les événements et parfois la supervision. Ils sont donc utiles lorsque l'objectif principal est d'orchestrer des workflows dans un moteur spécialisé.

Toutefois, construire une application web métier complète ne revient pas seulement à exécuter un workflow. Une application complète exige aussi une interface utilisateur adaptée, une architecture backend, une gestion des rôles, une persistance des données, des tests, des scripts de déploiement et une structure de code maintenable. Les travaux proches du développement web process-aware montrent que cette distinction est importante. Hemel et al. \cite{hemel2008} expliquent déjà que les langages de workflow ne sont pas toujours suffisants pour générer des applications completes, ce qui justifie l'extension d'un langage web par des abstractions de workflow.

Neto et al. \cite{neto2022} vont dans le même sens avec AKIP. Leur travail montre l’intérêt de plateformes permettant de développer des applications web process-aware, mais souligne aussi des limites récurrentes des solutions de Business Process Automation : verrouillage fournisseur, intégration complexe, interfaces limitées et dépendance technologique. Ces limites rejoignent directement le besoin de ProOWeb : produire une application dont le code source reste exploitable par les développeurs.

Ainsi, les BPMS ne sont pas inutiles ; ils répondent à un autre problème. Ils sont pertinents pour l’exécution et la supervision de workflows, mais ils ne garantissent pas toujours la génération d'une application web complète, modulaire et source-owned. ProOWeb se positionne donc comme une solution médiane : conserver l'apport de BPMN et de l'automatisation, sans imposer que l'application soit structurée autour d'un moteur externe central.

\mySection{Ingénierie dirigée par les modelés et génération de code}{}

L’ingénierie dirigée par les modèles fournit un cadre scientifique solide pour comprendre la transformation de modèles en artefacts logiciels. Dans cette approche, un modèle n'est pas seulement un support de documentation. Il peut être validé, transformé et utilisé pour produire du code, des configurations, des schémas, des tests ou des interfaces.

Les travaux de Syriani et al. \cite{syriani2018} montrent que la génération de code par templates est une technique largement utilisée dans les approches MDE. Elle permet de produire du code à partir de spécifications de haut niveau, mais elle soulève aussi des problèmes de qualité, de maintenance des templates et de traçabilité. Pour ProOWeb, cette observation est importante : il ne suffit pas de produire du code depuis BPMN ; il faut aussi gouverner la génération.

Siahaan et al. \cite{siahaan2024} étudient la transformation automatique entre artefacts de développement logiciel. Leur revue montre que les artefacts peuvent être reliés dans plusieurs directions : exigences vers modèles, modèles vers code, code vers tests, ou modèles vers documentation. Cette logique justifie l'introduction d'artefacts intermédiaires dans ProOWeb. Un modèle BPMN brut ne doit pas être transforme directement en fichiers applicatifs. Il doit d'abord être interprété et stabilisé dans des contrats intermédiaires.

Les travaux de Castillo et Davila \cite{castillo2024} sur le développement web dirige par les modèles confirment que la génération web est un domaine actif, mais héterogene. Les outils ne couvrent pas tous les mêmes artefacts. Certains se concentrent sur les interfaces, d'autres sur la navigation, d'autres sur les données ou la persistance. Cette héterogéneite montre qu'une approche comme ProOWeb doit definir clairement sa chaîne de transformation : BPMN enrichi, contrats, modules backend, composants frontend, tests et configuration.

\mySection{Génération, évolution et maîtrise du code}{}

La génération logicielle pose toujours la question du code produit. Si le code est illisible, difficile a modifier ou fortement lie à une plateforme, la productivité initiale peut être compensée par une dette de maintenance. Cette question est centrale pour ProOWeb, car le framework vise une application source-owned.

Les travaux de Hoyos et Restrepo-Calle \cite{hoyos2017} montrent que la génération automatique peut accélerer la production de prototypes web orientes processus. Uyanik et Sayar \cite{uyanik2023} proposent aussi une approche de gestion de processus web avec génération automatique de code. Ces travaux confirment que la génération d'applications web à partir de descriptions de processus est faisable et utile. Mais ils montrent aussi une limite : la génération se concentre souvent sur certains aspects, par exemple le prototype, les formulaires ou la gestion de flux, sans traiter complètement la question de l'évolution durable du code.

Dans les approches MDE, la maîtrise du code dépend fortement de la qualité des transformations, de la lisibilité des artefacts produits et de la possibilité d'intervenir manuellement. Dans les approches low-code, la question est souvent plus délicate : le code peut être masqué, difficilement exportable ou fortement dépendant de la plateforme. Pour ProOWeb, il faut donc séparer les artefacts gérés automatiquement des artefacts personnalises. Il faut aussi conserver une trace des fichiers générés, détecter les conflits et permettre une regénération sûre lorsque les modèles évoluent.

Cette exigence conduit à considérer la génération comme un processus continu, et non comme une opération ponctuelle. Une application générée doit pouvoir évoluer avec les processus métier. Le modéle BPMN peut changer, les règles d'accès peuvent être modifiées, les formulaires peuvent être enrichis, et les développeurs peuvent ajouter du code métier. Le framework doit donc prévoir une gouvernance de l'évolution.

\mySection{Low-code, no-code et model-driven engineering}{}

Les plateformes low-code et no-code visent à réduire l'effort de développement en s'appuyant sur des modèles visuels, des composants préconfigures et des mécanismes d'automatisation. Leur promesse est proche de celle des approches model-driven : augmenter le niveau d'abstraction et produire plus rapidement des applications.

Bucaioni et al. \cite{bucaioni2022} montrent que la modélisation occupe une place importante dans le low-code. Les plateformes low-code utilisent souvent des représentations graphiques ou déclaratives, même lorsqu'elles ne se réclament pas explicitement de l'ingénierie dirigée par les modèles. Di Ruscio et al. \cite{diruscio2022} analysent plus directement le lien entre low-code et MDE. Ils montrent que les deux approches partagent des objectifs communs, mais qu'elles diffèrent par leurs pratiques, leurs outils et leur rapport au code.

Martinez-Lasaca et al. \cite{martinez2026} proposent une approche model-driven et workflow-driven pour construire des plateformes low-code spécifiques à un domaine. Cette contribution est importante pour ProOWeb, car elle montre que le low-code peut être conçu avec des principes de modélisation rigoureux. Pankowska \cite{pankowska2024} rappelle cependant que le cycle de développement low-code donne souvent un rôle important aux utilisateurs finaux et modifie la frontière entre spécification et implémentation.

ProOWeb peut donc être situe à l'intersection de ces approches, mais sans se confondre avec elles. Il partage avec le low-code l'objectif d'accélération. Il partage avec le MDE l'importance des modèles, des transformations et des artefacts intermédiaires. Il s'en distingue par son orientation source-owned : l'application produite doit rester un projet logiciel que l'équipe peut lire, versionner, tester, modifier et déployer.

\mySection{Lacunes et positionnement de ProOWeb}{}

La revue met en évidence plusieurs lacunes. La première concerne l'insuffisance de BPMN seul. BPMN est un bon modèle de processus, mais il ne contient pas toutes les informations nécessaires à une application web complète. Les travaux sur BPMN vers IFML, sur les ontologies de domaine et sur la génération d'interfaces montrent qu'un enrichissement est nécessaire.

La deuxième lacune concerne la génération partielle. Certains travaux générent des interfaces, d'autres des prototypes, d'autres des workflows ou des artefacts techniques. Peu d'approches couvrent une chaîne integrée allant du modèle BPMN enrichi jusqu'au backend, au frontend, aux tests, à la sécurite, aux contrats et aux scripts de déploiement.

La troisième lacune concerne la maîtrise du code source. Les BPMS et plateformes low-code peuvent accélerer la construction d'applications, mais ils introduisent parfois une dépendance forte à un moteur ou à une plateforme. ProOWeb doit donc produire une application source-owned, c'est-a-dire un projet dont le code reste disponible, compréhensible et maintenable.

La quatrième lacune concerne l'évolution. Les travaux existants parlent souvent de génération, mais moins de regénération. Or les processus métier changent. Les développeurs modifient le code. Les besoins de sécurite, de formulaire ou de données évoluent. Une approche générative doit donc prévoir la préservation des personnalisations, la détection des conflits et la traçabilité entre modéles et fichiers.

La cinquième lacune concerne la modularité. Toutes les applications orientées processus n'ont pas les mêmes besoins. Certaines necessitent des notifications, d'autres une organisation hiérarchique, un audit, des tests end-to-end ou une intégration externe. Une architecture générative doit donc permettre d'activer ou non des capacités, ce qui justifie l'idée de modules ou de feature packs.

A partir de ces lacunes, ProOWeb peut être positionné comme un framework génératif orienté BPMN et applications web métier. Il ne s'agit pas d'un simple éditeur BPMN, ni d'un BPMS classique, ni d'une plateforme low-code fermée. ProOWeb se situe entre le développement manuel, les BPMS et les plateformes low-code. Il cherche à réduire la traduction manuelle entre modéle et code, tout en conservant une architecture logicielle claire et un code source maîtrisable.

\mySection{Conclusion}{}

Cette revue montre que plusieurs familles de travaux abordent le passage des modèles de processus vers des artefacts logiciels. BPMN fournit un langage de représentation métier. Les BPMS permettent d'exécuter les workflows. Les approches MDE et la génération par templates permettent de transformer des modéles en code ou en artefacts. Les travaux sur BPMN vers IFML montrent que les interfaces peuvent être dérivées de processus. Les plateformes low-code montrent l'intérêt industriel de l'automatisation.

Cependant, ces familles de travaux ne résolvent pas complètement le problème traité dans ce mémoire. Il reste difficile de générer une application web complète, maintenable, modulaire et source-owned à partir de modèles BPMN enrichis. Les difficultés concernent la couverture backend/frontend, l'intégration des données, la maîtrise du code, la regénération sûre et la gestion des personnalisations.

Ces constats justifient la conception de ProOWeb. Le framework vise à exploiter BPMN comme point d'entrée, à l'enrichir par des spécifications complémentaires, à produire des contrats intermédiaires, puis à générer une application web complète dont le code reste maîtrisable par l'équipe de développement. La suite du mémoire peut donc présenter l'approche d'ingéniérie, l'architecture du framework et les mécanismes de génération proposes.
